\documentclass[10pt,a4paper]{amsart}
\usepackage[margin=19mm]{geometry}
\usepackage{amsmath,amssymb,mathtools}
\usepackage[scr=rsfs,cal=euler]{mathalfa}
\usepackage{xfrac}
\usepackage[unicode,hidelinks,bookmarks=false]{hyperref}
\newcommand{\ie}{i.e.\ }
\newcommand{\cf}{cf.\ }
\newcommand{\resp}{respectively\ }
\newcommand{\Subsection}{\subsection}
\providecommand{\nobreakdash}{\nobreak\hbox{-}}
\setlength{\emergencystretch}{2em}
\multlinegap0pt
\usepackage[shortlabels]{enumitem}
\usepackage{amssymb}
\usepackage{xcolor}
\usepackage{algorithm}\usepackage{algpseudocode}
\usepackage{tikz}
\newtheorem{theorem}{Theorem}
\newcommand{\R}{\mathbb{R}}
\newcommand{\abs}[1]{\left\lvert #1 \right\rvert}
\newcommand{\argmin}{\operatorname*{arg\,min}}
\newcommand{\dist}{\operatorname{dist}}
\newcommand{\norm}[1]{\left\lVert #1 \right\rVert}
\newcommand{\pos}[1]{\left(#1\right)_{+}}
\title{Introduction to Exact Penalization for Mathematical Programming with Equilibrium Constraints}
\author{Louis Shuo Wang}
\date{Source: arXiv:2605.00387v1 (CC BY 4.0)}
\begin{document}
\maketitle

\noindent\emph{Source and licence.} Introduction to Exact Penalization for Mathematical Programming with Equilibrium Constraints. Version arXiv:2605.00387v1 (CC BY 4.0), distributed by arXiv under the Creative Commons Attribution 4.0 International licence, http://creativecommons.org/licenses/by/4.0/, as recorded in the arXiv licence metadata for this exact version and shown on the abstract page https://arxiv.org/abs/2605.00387v1. Full licence notice: \texttt{licenses/arxiv-2605.00387-CC-BY-4.0.txt}.\par\smallskip

\noindent\textit{Editorial scope.} Counted unit: the article's theorem thm:general-exact (source0.tex lines 261--277). The article's complete original proof of it is delivered with it (source0.tex lines 279--299, one \texttt{proof} environment of 21 lines). No mathematics is written, repaired or re-derived: the statement and the proof are the article's own lines, in the article's own notation, and no retained line is substituted or re-flowed.  Retained uncounted as the source-local context the proof needs: lines 167--174; lines 178--181; lines 233--243.  Nothing was omitted from any retained span, so the package carries no omission record.  No retained line contains a citation, so the wrapper carries no bibliography and no bibliography entry.  No retained line contains a figure environment, an image-inclusion command or drawing code, so no image is delivered and the package declares no asset dependency.  Editorial formatting only, in addition: an AMS \texttt{amsart} preamble that restates the article's own declarations and macros used by the retained lines, namely the environments \texttt{theorem} on separate counters, together with the standard packages those lines need.  The retained environments print in this document as Theorem 1 in order of appearance, whereas the article prints the same items with the numbering of its own full document.  The complete native source of the article as distributed in the arXiv e-print for this exact version is bundled unchanged as \texttt{sources/199556/source0.tex}.
\begin{equation}
\label{eq:nlp}
\begin{aligned}
    \min_x\quad & \theta(x)\\
    \text{s.t.}\quad & x\in X,\\
    & g(x)\le 0,\qquad h(x)=0,
\end{aligned}
\end{equation}

\begin{equation}
\label{eq:nlp-residual}
r(x):=\sum_{i=1}^p \pos{g_i(x)}+\sum_{j=1}^q \abs{h_j(x)}.
\end{equation}

\begin{theorem}[Łojasiewicz-type comparison inequality]
\label{thm:loja}
Let $S\subseteq \R^n$ be compact and subanalytic, and let $\phi,\psi:S\to\R$ be continuous and subanalytic. If
\[
\phi^{-1}(0)\subseteq \psi^{-1}(0),
\]
then there exist $\rho>0$ and an integer $N^\ast\ge 1$ such that
\[
\rho\,\abs{\psi(x)}^{N^\ast}\le \abs{\phi(x)}\qquad \forall x\in S.
\]
\end{theorem}

\begin{theorem}[General exact penalty theorem]
\label{thm:general-exact}
Consider \eqref{eq:nlp}. Assume:
\begin{enumerate}[label=(\roman*)]
    \item $X\subseteq \R^n$ is compact and subanalytic;
    \item $\theta$ is Lipschitz continuous on $X$;
    \item each $g_i$ and $h_j$ is continuous and subanalytic;
    \item the feasible set $W$ is nonempty.
\end{enumerate}
Then there exist $\alpha^\ast>0$ and an integer $N^\ast\ge 1$ such that for every $\alpha\ge \alpha^\ast$ and every $N\ge N^\ast$,
\[
\argmin_{x\in W}\theta(x)
=
\argmin_{x\in X}\Bigl(\theta(x)+\alpha r(x)^{1/N}\Bigr),
\]
where $r$ is the residual \eqref{eq:nlp-residual}.
\end{theorem}

\begin{proof}[Proof sketch]
Let
\[
\psi(x):=\dist(x,W).
\]
By subanalyticity and compactness, Theorem~\ref{thm:loja} implies that
\[
\rho\,\psi(x)\le r(x)^{1/N^\ast}\qquad \forall x\in X
\]
for some $\rho>0$. If $z\in W$ is a nearest feasible point to $x$, then Lipschitz continuity gives
\[
\theta(x)\ge \theta(z)-L\norm{x-z}=\theta(z)-L\psi(x).
\]
Hence
\[
\theta(x)+\alpha r(x)^{1/N^\ast}
\ge
\theta(z)+(\alpha\rho-L)\psi(x).
\]
Thus whenever $\alpha\rho>L$, the penalized objective cannot improve on a feasible comparator unless $\psi(x)=0$, i.e., unless $x\in W$. This proves exactness.
\end{proof}

\end{document}
